El área y el Hamiltoniano modular forman una pareja conceptualmente análoga a producto y cociente en el vacío.\par

El operador de área determina la ocupación total. El Hamiltoniano modular conserva la procedencia sectorial de esa ocupación.\par

En el borde aparecen dos familias de canales: \(90\) y \(120\). Cada una posee sus operadores de ocupación. El operador de área se construye a partir de la suma:\par

\[
\frac{\widehat{\mathcal A}_{\rm phys}}{\ell_P^2}
=
\gamma_{\rm BI}a_0
(N_{90}+N_{120}).
\]

El área cuenta el número total de excitaciones. Si una excitación procede del canal \(90\) o del canal \(120\), el área elemental es la misma. Desde el punto de vista geométrico, ambas contribuyen al mismo total.\par

El Hamiltoniano holográfico modular introduce una ponderación sectorial adicional:\par

\[
K_{\rm HHM}-cI
=
\Delta_4^{\rm mod}
(90N_{90}+120N_{120}).
\]

Cuenta las ocupaciones y pondera además cada excitación por el canal del que procede.\par

Por eso el área y el Hamiltoniano conmutan: ambos leen los mismos operadores de ocupación y contienen informaciones complementarias.\par

El espectro del operador de área es invariante respecto de la procedencia sectorial, mientras que el Hamiltoniano modular la distingue.\par

El operador de área determina el número total de excitaciones. El valor modular ponderado determina su distribución entre las dos genealogías. La conservación conjunta de ambos observables reconstruye exactamente la descomposición sectorial.\par

Ésta es una realización especialmente nítida de la idea holográfica. El borde publica un conjunto mínimo de observables que permite recuperar la información pertinente del interior.\par

La holografía consiste aquí en identificar qué par de observables fronterizos reconstruye exactamente qué datos internos.\par

El operador de área publica la magnitud geométrica y el Hamiltoniano modular publica la genealogía sectorial. La pareja de observables reconstruye la información completa de los sectores \(90/120\).\par

El estado reducido del borde convierte además ese Hamiltoniano en un verdadero generador modular:\par

\[
K_{\rm HHM}=-\log\rho_A.
\]

La purificación termocampo doble muestra que la entropía del estado reducido es exactamente entropía de entrelazamiento dentro del sistema duplicado declarado. Existe un estado puro ampliado cuya traza parcial produce el estado de borde y realiza explícitamente la identidad térmica.\par

Entonces aparece el déficit informacional orientado:\par

\[
I_{\rm or}
=
36\log3-S(\rho_A).
\]

Los \(36\) trits representarían la capacidad máxima si todos los estados fueran uniformes. La diferencia mide cuánta información orientada conserva la distribución real de pesos \(90/120\).\par

La entropía total cuantifica la incertidumbre. El déficit respecto del máximo cuantifica la estructura orientada residual.\par

La primera ley modular traduce esa estructura a geometría:\par

\[
\delta S
=
\beta_{90}\,
\delta\langle\mathcal A_{90}\rangle
+
\beta_{120}\,
\delta\langle\mathcal A_{120}\rangle,
\]

con\par

\[
\frac{\beta_{120}}{\beta_{90}}
=
\frac43.
\]

El cambio de entropía depende conjuntamente del cambio de área total y del sector de procedencia. Dos variaciones con el mismo incremento geométrico pueden tener distinta respuesta informacional si proceden de canales diferentes.\par

Esta dependencia establece la conservación geométrica de la memoria.\par

Para variaciones finitas aparece la entropía relativa como término positivo de corrección. La extensión exacta de la ley incorpora la diferencia entre la variación exacta y su aproximación lineal mediante una cantidad positiva.\par

Esta corrección positiva determina, fuera del régimen infinitesimal, el coste no lineal del alejamiento respecto del estado de referencia y cuantifica la contribución de memoria genealógica en la variación finita.\par

Finalmente, la relación\par

\[
E_P=k_B\Theta_{\rm clk}\log3
\]

y el cuanto areal\par

\[
\delta\mathcal A
=
\gamma_{\rm BI}a_0\ell_P^2
\]

permiten definir una tensión de conversión entre información energética y geometría:\par

\[
\mathcal T_{I/A}
=
\frac{E_P}{\delta\mathcal A}.
\]

Se expresa la razón exacta entre el cuanto energético informacional y el cuanto geométrico de área en la normalización declarada.\par

Por eso el Hamiltoniano holográfico modular es una pieza central dentro de un atlas de lectores complementarios. Su función específica consiste en conservar la procedencia que la geometría total integra en una suma.\par

El operador de área determina la ocupación total.\par
El Hamiltoniano modular determina la procedencia sectorial.\par
La entropía cuantifica la información publicada.\par
La entropía relativa cuantifica el coste de la transición genealógica.\par
